% =====================================================================
%  Thème 4 — Produit scalaire et orthogonalité — MOYEN — Corrigé
% =====================================================================
\documentclass[11pt,a4paper]{article}
\usepackage[utf8]{inputenc}
\usepackage[T1]{fontenc}
\usepackage{lmodern}
\usepackage[french]{babel}
\usepackage{amsmath,amssymb,mathtools}
\usepackage{xcolor}
\usepackage{enumitem}
\usepackage{fancyhdr}
\usepackage[a4paper,margin=2.1cm,headheight=26pt,headsep=10pt]{geometry}

\definecolor{primary}{RGB}{222,70,80}
\definecolor{primarydark}{RGB}{150,30,42}
\definecolor{excol}{RGB}{0,135,90}

\pagestyle{fancy}\fancyhf{}
\fancyhead[L]{\small\textcolor{primarydark}{\textbf{Fiche 2} — Calcul vectoriel}}
\fancyhead[R]{\small\textcolor{primarydark}{\textsc{Thème 4 — Moyen — Corrigé}}}
\fancyfoot[C]{\small\thepage}
\renewcommand{\headrulewidth}{0.8pt}

\newcommand{\vc}[1]{\overrightarrow{#1}}
\providecommand{\norm}[1]{\left\lVert #1\right\rVert}
\newcommand{\cor}[1]{\medskip\noindent{\color{primarydark}\bfseries\large Corrigé #1.}\par\smallskip}
\newcommand{\rep}{\textcolor{excol}{$\blacktriangleright$}\ }

\begin{document}
\begin{center}
{\LARGE\bfseries\color{primarydark}Produit scalaire et orthogonalité — Corrigé Moyen}\\[2pt]
{\color{primary}\rule{0.55\linewidth}{1.1pt}}
\end{center}
\vspace{2mm}

\cor{1}
\begin{enumerate}[label=\textbf{\alph*)},leftmargin=2em,itemsep=2pt]
  \item $\vec u\cdot\vec v=3\times2+y\times6=6+6y=0 \Rightarrow y=-1$.
  \item $x\times2+4\times(-1)=2x-4=0 \Rightarrow x=2$.
  \item $m\times m+2\times(-8)=m^2-16=0 \Rightarrow m=4 \text{ ou } m=-4$.
\end{enumerate}
\rep « Orthogonaux » se traduit par « produit scalaire nul » : on pose $xx'+yy'=0$ et on résout.

\cor{2}
\begin{enumerate}[label=\textbf{\alph*)},leftmargin=2em,itemsep=2pt]
  \item $\vc{AB}\,(3;1)$ et $\vc{AC}\,(1;4)$.
  \item $\vc{AB}\cdot\vc{AC}=3\times1+1\times4=3+4=7>0$ : l'angle $\widehat{BAC}$ est \textbf{aigu}.
\end{enumerate}

\cor{3}
\begin{enumerate}[label=\textbf{\alph*)},leftmargin=2em,itemsep=2pt]
  \item $d$ : $a=2,b=-3\Rightarrow$ vecteur directeur $\vec u\,(3;2)$.\\
        $d'$ : $a=3,b=2\Rightarrow$ vecteur directeur $\vec u'\,(-2;3)$.
  \item $\vec u\cdot\vec u'=3\times(-2)+2\times3=-6+6=0$ : les droites $d$ et $d'$ sont
        \textbf{perpendiculaires}.
\end{enumerate}
\rep On teste la perpendicularité de deux droites par le produit scalaire de leurs vecteurs directeurs.

\cor{4}
\begin{enumerate}[label=\textbf{\alph*)},leftmargin=2em,itemsep=2pt]
  \item $\vec u\cdot\vec v=2\times1+3\times(-1)=2-3=-1$ ; \quad
        $\vec w\cdot\vec v=4\times1+0\times(-1)=4$.
  \item $(2\vec u)\cdot\vec v=2(\vec u\cdot\vec v)=2\times(-1)=-2$.
        Direct : $2\vec u\,(4;6)$, $(4;6)\cdot(1;-1)=4-6=-2$. \checkmark\\
        $(\vec u+\vec w)\cdot\vec v=\vec u\cdot\vec v+\vec w\cdot\vec v=-1+4=3$.
        Direct : $\vec u+\vec w\,(6;3)$, $(6;3)\cdot(1;-1)=6-3=3$. \checkmark
\end{enumerate}
\rep La linéarité permet de « distribuer » le produit scalaire, comme une multiplication.

\cor{5}
\begin{enumerate}[label=\textbf{\alph*)},leftmargin=2em,itemsep=2pt]
  \item $\vc{AB}\,(2;0)$, $\vc{AD}\,(0;2)$ : $\vc{AB}\cdot\vc{AD}=2\times0+0\times2=0$. Les côtés
        $[AB]$ et $[AD]$ sont \textbf{perpendiculaires} (angle droit du carré).
  \item $\vc{AC}\,(2;2)$ : $\vc{AB}\cdot\vc{AC}=2\times2+0\times2=4$.
  \item $\vc{BD}\,(0-2;2-0)=(-2;2)$ : $\vc{AC}\cdot\vc{BD}=2\times(-2)+2\times2=-4+4=0$.
        Les diagonales $[AC]$ et $[BD]$ du carré sont \textbf{perpendiculaires}.
\end{enumerate}
\rep Dans un carré, côtés adjacents et diagonales sont perpendiculaires : le produit scalaire le
confirme numériquement.

\end{document}
